\documentclass[11pt,letterpaper]{article}
\usepackage[T1]{fontenc}
\usepackage[utf8]{inputenc}
\usepackage{lmodern}
\usepackage[margin=1.02in,headheight=14pt]{geometry}
\usepackage{amsmath,amssymb,amsthm,mathtools}
\usepackage{microtype}
\usepackage{booktabs,longtable,array}
\usepackage{xcolor}
\usepackage{enumitem}
\usepackage{listings}
\usepackage{fancyhdr}
\usepackage{xurl}
\usepackage[colorlinks=true,linkcolor=blue!45!black,citecolor=blue!45!black,
  urlcolor=blue!45!black,pdftitle={Cayley--Hamilton in Intuitionistic Primitive Recursive Arithmetic},
  pdfauthor={Research report prepared for Vladimir Reshetnikov}]{hyperref}
\newtheorem{theorem}{Theorem}[section]
\newtheorem{proposition}[theorem]{Proposition}
\newtheorem{lemma}[theorem]{Lemma}
\newtheorem{corollary}[theorem]{Corollary}
\theoremstyle{remark}
\newtheorem{remark}[theorem]{Remark}
\setlist{itemsep=3pt,topsep=5pt,leftmargin=*}
\lstset{basicstyle=\small\ttfamily,columns=fullflexible,
  keepspaces=true,breaklines=true,frame=single,rulecolor=\color{black!20},
  backgroundcolor=\color{black!2},showstringspaces=false}
\pagestyle{fancy}
\fancyhf{}
\fancyhead[L]{\small Cayley--Hamilton in iPRA}
\fancyhead[R]{\small Finite proofs and formalization}
\fancyfoot[C]{\thepage}
\setlength{\parskip}{3pt}
\setlength{\emergencystretch}{2em}
\setcounter{tocdepth}{2}
\numberwithin{equation}{section}

\begin{document}
\hypersetup{pageanchor=false}
\begin{titlepage}
\vspace*{0.6in}
{\Large\sffamily\color{blue!45!black}CONSTRUCTIVE ALGEBRA}\par
\vspace{0.35in}
{\huge\bfseries Cayley--Hamilton in\par
 Intuitionistic Primitive\par Recursive Arithmetic}\par
\vspace{0.25in}
{\Large Finite proofs, effective applications,\par
 and an explicit Lean formalization route}\par
\vspace{0.55in}
{\large Prepared for Vladimir Reshetnikov}\par
{\large 15 November 2026}\par
\vspace{0.5in}
\begin{abstract}
We prove a uniform arithmetized Cayley--Hamilton theorem using finite
determinant expansion, the adjugate identity, and coefficient telescoping.
Explicit coding reduces the theorem to a primitive recursive numerical
equation with quantifier-free induction invariants. The development
covers 27 result families: polynomial adjugates, monic reduction, unit
inversion, effective minimal polynomials, Newton and Bell identities,
Hessenberg formulas, Faddeev--LeVerrier, finite formal logarithms,
module and integrality consequences, interpolation, quantitative matrix
functions, and specified quaternionic extensions. Every use of division,
spectral data, coordinates, or analytic convergence is made explicit.
The accompanying Lean sources define the numerical checker, the exact
object-language target, and ordinary Mathlib comparison statements.
They are an uncompiled outline, not a completed Lean verification.
Exact executable checks and pinned source provenance accompany the proof.
\end{abstract}
\vfill
\noindent\textbf{Status.} Developed human proofs with an iPRA induction audit;
executable finite checks passed; complete Lean proof construction remains.
\end{titlepage}
\hypersetup{pageanchor=true}

\tableofcontents
\clearpage
\section*{Coverage and qualifications at a glance}
\addcontentsline{toc}{section}{Coverage and qualifications at a glance}

The starting inventory is the mathematical content of the specified
Wikipedia article, revision \texttt{1375228092}~\cite{wikipedia}.
The proofs below are developed afresh as finite algebra and finite
arithmetic. Classical background and related formulations can be found
in the Stacks Project and Conrad's notes~\cite{stacks,conrad}; their
abstract existence arguments are not imported as iPRA lemmas.

In the table, \textbf{U} means a universal finite ring identity, interpreted
arithmetically as an integer-polynomial identity; \textbf{C} means a
finite construction from the stated certificates in a fixed effective
presentation; and \textbf{Q} means a quantitative rational-approximation
version with a supplied or explicitly proved modulus. Every row refers
to a human proof in this manuscript. None denotes a completed Lean
kernel proof.

\small
\begin{longtable}{@{}p{0.07\textwidth}p{0.39\textwidth}p{0.45\textwidth}@{}}
\toprule
ID & Result family & Exact scope \\
\midrule
\endfirsthead
\toprule ID & Result family & Exact scope \\ \midrule
\endhead
\bottomrule
\endfoot
A1 & Determinant toolkit & U: signs, alternation, cofactors, adjugates, multiplicativity. \\
A2 & Characteristic coefficients & U: leading and constant coefficients, principal minors, trace coefficient. \\
A3 & Cayley--Hamilton & U: all dimensions, arbitrary commutative rings, including the zero ring. \\
A4 & Polynomial adjugate & U: every adjugate coefficient is an explicit polynomial in the matrix. \\
A5 & Evaluation and power reduction & U/C: restricted noncommutative evaluation; monic remainder and explicit power recurrence. \\
A6 & Inversion and negative powers & C: a supplied inverse of the determinant; exact numerical examples. \\
A7 & Minimal polynomials & C: bounded elimination over a fixed effective field; divisibility of the characteristic polynomial. \\
A8 & Companion matrices & U: prescribed characteristic polynomial and exact annihilator ideal. \\
A9 & Eigenvectors and roots & C: equivalence over an effective field; ring consequences with a unit component. \\
A10 & Similarity and split spectra & U/C: supplied inverses, triangularization or factorization; coefficient and power-trace formulas. \\
T1 & Jacobi and adjugate traces & U: finite formal differentiation. \\
T2 & Newton identities & U: all degrees, without division or eigenvalues. \\
T3 & Bell-polynomial formulas & U: factorial-multiplied identities; division only with justified units or exact integer quotients. \\
T4 & Hessenberg determinant formula & U: determinant equals $k!e_k$; explicit formulas through order four. \\
T5 & Faddeev--LeVerrier & C: correctly initialized; supplied units or proved exact integer division. \\
T6 & Exponential/logarithm and partitions & C: every finite formal precision, with integer denominators invertible. \\
M1 & Finitely generated modules & C: supplied generator images and coordinates; no basis required. \\
M2 & Ideal coefficients and Nakayama & C: finite membership and unit certificates. \\
M3 & Closure of integral elements & C: monic input equations give a degree-$\prod_jd_j$ output equation. \\
M4 & Algebraic-number minimal polynomials & C: finite multiplication table; integral coefficients for integral elements. \\
I1 & Lagrange and Sylvester formulas & C: supplied roots with unit differences; spectral idempotents. \\
I2 & Hermite interpolation & C: repeated-root jets, Hasse derivatives, explicit truncated inverses. \\
F1 & Polynomial matrix-function reduction & Q: primitive recursive rational coefficients and weighted finite-tail bound. \\
F2 & Exponential, sine, cosine & Q: rational matrix and parameter inputs, explicit factorial cutoff. \\
F3 & Rotation formulas & Q: planar, Pauli, and Rodrigues identities from supplied polynomial relations. \\
Q1 & Quaternion and split-quaternion matrices & C: faithful representations give degree $2n$ annihilators; also degree $4n$ regular representation. \\
Q2 & One octonion & U/C: quadratic norm identity, with rational coordinate interpretation. \\
\end{longtable}
\normalsize

The unrestricted analytic functional calculus is not claimed as a literal
first-order iPRA theorem. Neither are a root-finding or Jordan-form
existence theorem for arbitrary fields, density of diagonalizable complex
matrices, or the modified matrix-octonion theorem. The article's broader
Lie-group references and its exterior-algebra and combinatorial proof
methods are not developed here. The main theorem already has two
complete finite proofs; those other methods are unnecessary dependencies.

Several source formulations require correction. Over a ring, an inverse
requires a unit determinant. Trace expressions with factorial
denominators must be qualified. The stated partition product has an
extra sign, and the displayed Faddeev--LeVerrier recursion lacks its
first valid initialization. A nonpolynomial analytic function cannot be
written $qp+r$ with both $q,r$ polynomials. These points are proved or
diagnosed in the appropriate sections rather than used as premises.


\clearpage
\section{What is proved, and what ``in iPRA'' means}
\label{sec:foundations}

The main result is a uniform arithmetic form of the Cayley--Hamilton theorem.
It is accompanied by a finite constructive development of its principal
algebraic consequences. Throughout, a ring is associative, commutative, and
unital unless explicitly stated otherwise. We allow the zero ring. Thus a
polynomial displayed with leading coefficient $1$ is monic in the formal
coefficient sense; its actual degree is $n$ only when $1\ne0$.

\begin{theorem}[Uniform arithmetic Cayley--Hamilton]
\label{thm:arithmetic}
There is a primitive recursive function $\mathrm{CHCheck}:\mathbb N\to
\{0,1\}$, defined explicitly below by integer polynomial arithmetic, for which
intuitionistic primitive recursive arithmetic proves
\[
                 \forall n\;\mathrm{CHCheck}(n)=1.
\]
The assertion says that every coefficient in every entry of
$\det(tI_n-X)\big|_{t=X}$ is zero, where $X=(x_{ij})$ is the generic matrix
over $\mathbb Z[x_{00},\ldots,x_{n-1,n-1}]$. Consequently, specializing the
variables gives $\chi_A(A)=0$ over every commutative ring. For each fixed
primitive recursive presentation of a ring whose ring laws are provable in
iPRA, this specialization is itself an iPRA derivation.
\end{theorem}

The human proof of this theorem is given in Sections~\ref{sec:determinants}
and~\ref{sec:main}, with its induction accounting below. The accompanying
Lean files are a specification and a route to a completed formalization;
they are not a completed Lean kernel verification of this arithmetic
theorem. The executable checks are independent finite corroboration, not
a substitute for the uniform proof.

There are three distinct levels of assertion in this report.
\begin{enumerate}
\item \emph{Arithmetic identities}: equations between specified primitive
recursive functions on natural-number codes. These are the literal iPRA
statements.
\item \emph{Algebraic interpretations}: theorems about arbitrary rings,
modules, or supplied field data obtained from the finite identities. Such
objects are mathematical interpretations, not new sorts silently added to
first-order arithmetic.
\item \emph{Effective analytic versions}: statements about rational
approximations with explicit primitive recursive moduli. These are reduced
to finite inequalities at any requested precision, not quantified
existence claims about arbitrary real numbers or arbitrary analytic
functions.
\end{enumerate}
This distinction is necessary: a short theorem in ordinary Lean with
\texttt{[CommRing R]} does not alone establish an object-language iPRA proof.

\subsection{The base theory and its permitted induction}

We use intuitionistic first-order logic with equality, symbols for the
primitive recursive functions, their defining equations, and induction
for quantifier-free formulas with parameters. A primitive recursive
definition is built from zero, successor, projections, composition, and
primitive recursion. Abbreviations for newly defined primitive recursive
functions are conservative definitional extensions. This is the
presentation underlying the repository interface inspected in
Section~\ref{sec:lean}.

The quantifier-free induction rule has the form
\[
 \frac{\varphi(0,\mathbf a)\qquad
       \varphi(k,\mathbf a)\longrightarrow\varphi(k+1,\mathbf a)}
      {\varphi(k,\mathbf a)},
 \qquad \varphi\text{ quantifier-free}.
\]
Universal closure is taken after the derivation, not inserted into its
induction formula. Equality of natural numbers and every primitive
recursive Boolean test are decidable by their defining recursions. Thus
finite case distinctions and bounded search are legitimate intuitionistically.
No step requires excluded middle for arbitrary formulas, Markov's
principle, choice, an infinite basis construction, or a completeness axiom.

If $P(i,\mathbf a)$ is a decidable primitive recursive predicate, the bounded
assertion $\forall i<b\,P(i,\mathbf a)$ can be replaced by the equation
\[
 \operatorname{all}(b,P,\mathbf a)=1,\qquad
 \operatorname{all}(0,P,\mathbf a)=1,\quad
 \operatorname{all}(k+1,P,\mathbf a)
 =\operatorname{all}(k,P,\mathbf a)\,\mathbf1_{P(k,\mathbf a)}.
\]
Here $P$ is a fixed definition, not a quantified function argument.
An elementary quantifier-free induction proves that the product is $1$
exactly when each test is $1$. Nested bounded assertions are flattened in
the same way. This device will express statements about all entries,
coefficients, permutations, and intermediate stages of a finite calculation.

\subsection{Explicit coding conventions}

One possible list coding uses
\[
 \operatorname{nil}=0,\qquad
 \operatorname{cons}(a,b)=2^a(2b+1).
\]
For a positive code, find its exponent of $2$ by searching up to the code;
the remaining odd factor determines the tail. The tail is smaller than
the code. Repeated decoding for at most the original code many steps
therefore suffices to define length, indexing with default $0$, append,
and bounded updates by primitive recursion. Efficient coding is
irrelevant to proof-theoretic strength; this deliberately simple coding
also shows that every advertised bound is a natural-number bound.

Integers are encoded by $2a$ for $a\ge0$ and $2a+1$ for $-(a+1)$.
Signed addition, multiplication, negation, and comparison are primitive
recursive, and their ring laws have the familiar quantifier-free proofs
from the laws of natural-number arithmetic. A rational is a signed
integer numerator and a positive denominator, optionally reduced by the
Euclidean algorithm. That algorithm can be run for a bound supplied by
its initial arguments, with a stable state after termination. Comparison
and equality use cross multiplication. Invalid input codes receive a
fixed default so that all underlying functions are total.

An $n\times m$ matrix is a list of $nm$ scalar codes, indexed in row-major
order. A univariate polynomial is a coefficient list, interpreted with
zero beyond its length. Matrix multiplication, polynomial convolution,
trace, and powers are bounded sums and products or primitive recursions.
Polynomial coefficients need not be trimmed: degree \emph{bounds}, rather
than a choice of the last nonzero coefficient, suffice throughout the core.

For a completely literal generic identity, use multivariate polynomials
over the coded integers. With $r$ variables and a coordinate exponent
bound $d$, enumerate exponent vectors $\alpha\in\{0,\ldots,d\}^{r}$
by the base-$(d+1)$ number
\[
                    [\alpha]=\sum_{j<r}\alpha_j(d+1)^j.
\]
Store the coefficient array of length $(d+1)^r$. Addition pads to a common
bound. Multiplication increases the bound to the sum of the input bounds
and uses the finite convolution
\[
 (fg)_\alpha=\sum_{\beta,\gamma:\,\beta+\gamma=\alpha}
                       f_\beta g_\gamma.
\]
The sum is implemented over all bounded vector pairs with a Boolean
filter. Thus no unproved normal-form decision procedure is assumed.
Commutativity, associativity, and distributivity follow by finite
reindexing of these convolution sums. At a fixed degree bound, polynomial
equality is equality of two finite arrays.

For dimension $n$, set $r=n^2$ and form the generic matrix $X$. Compute
$c_0,\ldots,c_n$ from the permutation expansion of $\det(tI-X)$, then
compute $\sum_{j=0}^{n}c_jX^j$. The coefficient $c_j$ is supported only on
monomials of total degree $n-j$, and each entry of $X^j$ is supported
only in total degree $j$. Product expansion and induction on $j$ prove
these support assertions, including zero entries. Thus every residual
entry is supported in total degree $n$. No coefficient is discarded by
the coordinate exponent cap $n$, and comparing
the $n^2(n+1)^{n^2}$ integer coefficients suffices. For $n=0$ this is an
empty conjunction, equal to $1$. This algorithm is the promised
$\mathrm{CHCheck}$.

\subsection{Why primitive recursiveness alone is insufficient}
\label{sec:induction}

An algorithm being primitive recursive does not imply that every true
claim about it is provable in iPRA. The proof obligations here are
specific finite loop invariants, not an appeal to that invalid principle.
The following bookkeeping explains how the ordinary notation in the
proofs abbreviates quantifier-free arithmetic reasoning.

Fix an ambient dimension $n$ and the inputs to a claimed identity. All
index sets have explicit bounds: matrix indices below $n$, candidate
permutations among the $n^n$ maps $\{0,\ldots,n-1\}\to\{0,\ldots,n-1\}$,
subsets among $2^n$ bit masks, and polynomial coefficient arrays bounded
by the displayed exponent caps. Use $0^0=1$ for the unique empty map.
The following invariants suffice.
\begin{enumerate}
\item A finite sum after $k$ steps equals the sum of its first $k$ summands;
reindexing follows by inserting each value of a bounded permutation
exactly once. Products have the corresponding invariant.
\item In a sign-reversing involution, the sum over the first $k$ pairs is
zero. Pairs are indexed by the smaller member of their two-element orbit,
selected by a decidable numerical comparison. This proves cancellation
even in characteristic $2$, without cancelling a factor $2$.
\item In convolution, the partial sum accounts for exactly the first $k$
index pairs or triples. Associativity is a reindexing of the same bounded
triple sum, so its packed coefficient invariant is quantifier-free.
\item The adjugate entry identity is the cofactor expansion with a row or
column replaced. Its bookkeeping is the bounded permutation deletion
map described in Section~\ref{sec:determinants}.
\item In the main proof, after processing coefficient $k<n$, the invariant
is exactly
\[
                \sum_{j=0}^{k} c_jX^j=-X^{k+1}B_k.
\]
It is a finite array equality. No induction on the formula ``all matrices
of size $k$ satisfy Cayley--Hamilton'' is used.
\end{enumerate}
These are equations between fixed primitive recursive functions, or
bounded conjunctions of such equations. If a lower-dimensional identity
is needed in a dimension induction, keep $n$ fixed and pack all its
restrictions of sizes at most the current stage into one finite array
invariant. Alternatively, the determinant identities below are derived
directly at the given dimension by finite reindexing.

Integer sizes and list codes do not create an implicit unbounded search.
For the generic calculations, expand a computation without collecting
terms: addition adds term counts and multiplication multiplies them;
the sum of absolute integer coefficients obeys the same bounds.
Primitive recursion through the given computation produces an explicit
bound on every coefficient and code. For other fixed primitive
recursive scalar operations, monotone majorants are constructed along
their definitions. The finite computation is then indexed within those
bounds. This is the standard direct compilation of each stated
invariant into quantifier-free induction, not a uniform total interpreter
for all primitive recursive programs.

Universal polynomial equality can subsequently be evaluated in any
commutative ring: evaluate each finite monomial and sum, using the ring
laws to preserve addition and multiplication. This external
specialization uses no decision of equality in that ring. To make it an
internal iPRA theorem, fix a presentation with primitive recursive
operations and iPRA proofs of its ring laws. Such a presentation uses
either canonical scalar codes with numerical equality, or a primitive
recursive equivalence relation together with proofs of equivalence,
congruence, and the ring laws modulo that relation. Unreduced rational
pairs use the latter option, with cross-multiplication equality.
One must not assert that an
arbitrary program code describes such a ring, or that iPRA decides this.


\section{A finite determinant toolkit}
\label{sec:determinants}

Write $[n]=\{0,\ldots,n-1\}$. A permutation is a list containing each
element of $[n]$ once, and its sign is $(-1)^{\operatorname{inv}(\sigma)}$,
where an inversion is a pair $i<j$ with $\sigma(i)>\sigma(j)$. Define
\[
 \det M=\sum_{\sigma\in S_n}\operatorname{sgn}(\sigma)
                              \prod_{i<n}m_{i,\sigma(i)}.
 \tag{D}
\]
In arithmetic the sum can run through $n^n$ candidate lists and assign
zero to every list that is not a permutation. The determinant of the
empty matrix is $1$.

\begin{lemma}[Permutation signs]
Swapping two positions in a permutation reverses its sign. Moreover,
$\operatorname{sgn}(\sigma\tau)=\operatorname{sgn}(\sigma)
\operatorname{sgn}(\tau)$ and
$\operatorname{sgn}(\sigma^{-1})=\operatorname{sgn}(\sigma)$.
\end{lemma}
\begin{proof}
An adjacent interchange changes the inversion count by exactly $1$ in
absolute value. Swapping positions $i<j$ is a product of
$2(j-i)-1$ adjacent interchanges, hence changes parity. Bubble sorting a
list uses at most $n(n-1)/2$ adjacent interchanges and reduces its
inversion count to zero. Applying the same adjacent interchanges to
the positions in any other permutation multiplies its sign by the same
factor. This proves the product formula, then the inverse formula by
applying it to $\sigma\sigma^{-1}$. Every procedure and parity invariant
here is bounded in $n$.
\end{proof}

\begin{lemma}[Multilinearity, alternation, and transpose]
The determinant is additive and scalar-linear in each individual row or
column. It vanishes when two rows, or two columns, are equal.
Also $\det(M^{\mathsf T})=\det M$ and $\det I_n=1$.
\end{lemma}
\begin{proof}
Each summand of (D) contains one entry from a prescribed row and one
from a prescribed column, so distributivity proves multilinearity.
If rows $i$ and $j$ agree, pair $\sigma$ with the permutation obtained by
swapping its values at $i,j$. The products agree and the signs are
opposite. There are no fixed points. Sum once over the smaller member
of each pair to get zero. This argument never infers $z=0$ from $2z=0$.
The column proof is the same, or follows by transposition. Reindexing
by $\sigma^{-1}$ proves transposition invariance. For the identity matrix
only the identity permutation contributes, and its contribution is $1$.
\end{proof}

For $n\ge1$, write $M^{\widehat i,\widehat j}$ for the matrix obtained by
deleting row $i$ and column $j$, retaining their original orders. Define
\[
 C_{ij}=(-1)^{i+j}\det M^{\widehat i,\widehat j},\qquad
 (\operatorname{adj}M)_{ji}=C_{ij}.
\]

\begin{lemma}[Cofactor expansion and the two adjugate identities]
\label{lem:adjugate}
For every row $i$ and column $j$,
\[
 \det M=\sum_{k<n}m_{ik}C_{ik}
       =\sum_{k<n}C_{kj}m_{kj},
 \qquad
 M\operatorname{adj}M=\operatorname{adj}M\,M=(\det M)I_n.
\]
\end{lemma}
\begin{proof}
Partition the permutations according to the value $\sigma(i)=k$.
Delete the position $i$ and the value $k$, then relabel the remaining
ordered sets by $[n-1]$. This is a bijection with $S_{n-1}$; its inverse
inserts the missing position and value. Moving row $i$ to the first
position uses $i$ adjacent interchanges, and moving column $k$ to the
first position uses $k$; the surviving sign is therefore
$(-1)^{i+k}$ times the smaller permutation sign. The corresponding
products agree after extracting $m_{ik}$. This proves row expansion;
transpose gives column expansion.

The $(i,j)$ entry of $M\operatorname{adj}M$ is
$\sum_k m_{ik}C_{jk}$. It is the expansion along row $j$ of the matrix
obtained by replacing row $j$ of $M$ with row $i$. If $i=j$, it is
$\det M$; otherwise the replacement matrix has two equal rows and its
determinant is zero. Column replacement proves the other identity.
For $n=1$ the deleted determinant is $1$, so the proof includes that case.
\end{proof}

\begin{lemma}[Multiplicativity]
\label{lem:detmul}
For square matrices over a commutative ring,
$\det(AB)=\det A\det B$.
\end{lemma}
\begin{proof}
Column $j$ of $AB$ is $\sum_i b_{ij}A_i$, where $A_i$ is column $i$ of
$A$. Repeated multilinearity expands its determinant as a sum indexed
by maps $f:[n]\to[n]$, with term
\[
                 \left(\prod_j b_{f(j),j}\right)
                        \det(A_{f(0)},\ldots,A_{f(n-1)}).
\]
A noninjective $f$ repeats a column and contributes zero. Since its
domain and codomain have the same finite size, an injective $f$ is a
permutation: the occupancy counts $\nu_i=|f^{-1}(i)|$ all belong to
$\{0,1\}$ and sum to $n$, so a missing value would make their sum at
most $n-1$. Reordering its columns gives
$\operatorname{sgn}(f)\det A$. Factoring out $\det A$ leaves the column
form of (D) for $B$, proving the formula. The expansion has $n^n$ terms,
and is valid in the zero ring as well.
\end{proof}

\begin{corollary}[Change of scalars and similarity]
Every ring homomorphism preserves determinants and characteristic
polynomials. If $P,Q$ are supplied with $PQ=QP=I$, then
$\chi_{PAQ}=\chi_A$ and $f(PAQ)=P f(A)Q$ for every scalar polynomial $f$.
\end{corollary}
\begin{proof}
A homomorphism preserves the finite sum and products in (D).
The equality $tI-PAQ=P(tI-A)Q$ holds in the polynomial matrix ring, and
$\det P\det Q=1$ by Lemma~\ref{lem:detmul}. This proves the polynomial
identity. Induction gives $(PAQ)^k=PA^kQ$ for all $k$, including $k=0$;
summing proves the evaluation identity.
\end{proof}

\section{Cayley--Hamilton without division or spectral theory}
\label{sec:main}

Let $n\ge1$ and $A\in M_n(R)$. Work with finite polynomial lists in a
central indeterminate $t$. Set
\[
 \chi_A(t)=\det(tI-A)=\sum_{j=0}^n c_jt^j,
 \qquad B(t)=\operatorname{adj}(tI-A)=\sum_{j=0}^{n-1}B_jt^j.
\]
The degree bound for $B$ follows because every cofactor is a product
sum with $n-1$ linear factors. In $\chi_A$ a term of degree $n$ must
select $t$ from all $n$ rows. Only the identity permutation permits this,
and the coefficient is $1$. Thus $c_n=1$. At $t=0$,
$c_0=\det(-A)=(-1)^n\det A$.

\begin{theorem}[Coefficient telescoping]
\label{thm:CH}
For every square matrix over every commutative ring,
\[
       \boxed{\;A^n+c_{n-1}A^{n-1}+\cdots+c_1A+c_0I=0.\;}
\]
\end{theorem}
\begin{proof}
Lemma~\ref{lem:adjugate}, applied over the polynomial ring, gives
\[
                    (tI-A)B(t)=\chi_A(t)I.
\]
Comparing the finite coefficient lists gives
\begin{align*}
 B_{n-1}&=I,\\
 B_{j-1}-AB_j&=c_jI &&(1\le j\le n-1),\\
 -AB_0&=c_0I.
\end{align*}
Multiply the equation indexed by $j$ on the left by $A^j$ and add,
including the leading equation multiplied by $A^n$. The left side is
\[
 -AB_0+\sum_{j=1}^{n-1}
       \bigl(A^jB_{j-1}-A^{j+1}B_j\bigr)+A^nB_{n-1}=0;
\]
every summand cancels with its neighbor. The right side is precisely
$\sum_{j=0}^n c_jA^j$. In detail, the partial-sum invariant is
$\sum_{j=0}^k c_jA^j=-A^{k+1}B_k$ for $k<n$, followed by
$B_{n-1}=I$. This is the quantifier-free invariant identified in
Section~\ref{sec:induction}; it proves Theorem~\ref{thm:arithmetic} after
generic polynomial coding. No assertion about eigenvalues is used.

For $n=0$, the matrix algebra has one element, its zero and identity
coincide, and $\chi_A(t)=1$. Its evaluation is the empty identity matrix,
which is also zero. Thus the theorem includes dimension zero.
\end{proof}

\begin{theorem}[Adjugate as a polynomial]
\label{thm:adjpoly}
For $n\ge1$ and $0\le j<n$,
\[
 B_j=\sum_{k=j+1}^n c_k A^{k-1-j}.
\]
In particular,
\[
 \operatorname{adj}(-A)=A^{n-1}+c_{n-1}A^{n-2}+\cdots+c_1I,
 \qquad
 \operatorname{adj}A=(-1)^{n-1}
              \sum_{k=1}^{n}c_kA^{k-1}.
\]
\end{theorem}
\begin{proof}
Start with $B_{n-1}=c_nI=I$. If the formula holds for $B_j$, the
coefficient recurrence gives $B_{j-1}=AB_j+c_jI$, which is exactly the
next formula. This is a descending bounded recursion, or an ordinary
induction on $n-1-j$. The constant coefficient of the adjugate is
$\operatorname{adj}(-A)$, and each of its minors has $n-1$ negated
factors, giving the second identity. In particular every $B_j$ commutes
with $A$; this is a proved conclusion, not an assumption used above.
\end{proof}

\subsection{A second proof and the evaluation pitfall}

For a polynomial $F(t)=\sum F_it^i$ with coefficients in an associative
ring, define right evaluation by $F(A)=\sum F_iA^i$. It always preserves
addition, but generally does not preserve multiplication. If every
coefficient of $G$ commutes with $A$, however,
\[
 F(A)G(A)=\sum_{i,j}F_iA^iG_jA^j
         =\sum_{i,j}F_iG_jA^{i+j}=(FG)(A).
\]
This proves the required restricted substitution rule by a finite
calculation. In particular, evaluation is a ring homomorphism on
polynomials whose coefficients commute with $A$, even if those
coefficients do not commute with one another.

The two adjugate identities imply
$(tI-A)B(t)=B(t)(tI-A)$, hence $AB_j=B_jA$ by coefficient comparison.
The restricted substitution rule now legitimately evaluates
$(tI-A)B(t)=\chi_A(t)I$ at $A$, giving $0=\chi_A(A)$. This is a second,
fully finite proof. The stronger formula of Theorem~\ref{thm:adjpoly}
places these coefficients in the commutative algebra generated by $A$.

There is a concrete obstruction to unrestricted substitution. Take
\[
 A=\begin{pmatrix}0&1\\0&0\end{pmatrix},\qquad
 D=\begin{pmatrix}0&0\\1&0\end{pmatrix},\qquad F(t)=tI,\quad G(t)=D.
\]
Then $(FG)(A)=DA$, whereas $F(A)G(A)=AD$; the two products are different
diagonal idempotents. Thus a proof that treats evaluation on all
matrix-coefficient polynomials as a homomorphism is invalid.
Nor can one replace the scalar variable in $\det(tI-A)$ by a matrix
inside the entries and call the resulting expression $\det(A-A)$:
the types and the surrounding matrix sizes have changed.

An additional arithmetic diagnostic is provided by the permanent. For
$H=\left(\begin{smallmatrix}1&2\\3&4\end{smallmatrix}\right)$,
\[
 \operatorname{perm}(tI-H)=t^2-5t+10,
 \qquad H^2-5H+10I=12I\ne0.
\]
The same putative substitution proof would incorrectly annihilate this
permanent polynomial. Our proof specifically uses the sign cancellation
in the determinant.

\subsection{Inverse matrices, including the exact hypothesis}

\begin{theorem}[Explicit inverse from a unit determinant]
\label{thm:inverse}
For $n\ge1$, if $u\det A=1$ is supplied, then
\[
 A^{-1}=u\operatorname{adj}A
       =u(-1)^{n-1}
                (A^{n-1}+c_{n-1}A^{n-2}+\cdots+c_1I)
\]
is a two-sided inverse. Conversely a supplied inverse $V$ gives the
unit witness $\det V$ for $\det A$.
\end{theorem}
\begin{proof}
Multiply each adjugate identity by the scalar $u$ for the first claim.
The explicit polynomial expression follows from
Theorem~\ref{thm:adjpoly}. Conversely $AV=I$ gives
$\det A\det V=\det I=1$ by determinant multiplicativity. All witnesses
are explicit finite expressions.
\end{proof}
Over a field a nonzero determinant has an inverse, so this recovers the
usual statement. Over $\mathbb Z$, the matrix $(2)$ has nonzero
determinant and no inverse over $\mathbb Z$; nonzeroness is not the
correct ring hypothesis. The theorem remains valid in the zero ring,
where the required equations have their usual algebraic meanings.

\subsection{The displayed small examples}

For $A=(a)$ one has $\chi_A(t)=t-a$ and $A-aI=0$. For
$A=\left(\begin{smallmatrix}a&b\\c&d\end{smallmatrix}\right)$,
\[
 \chi_A(t)=t^2-(a+d)t+(ad-bc),\qquad
 A^2=\begin{pmatrix}a^2+bc&ab+bd\\ca+dc&cb+d^2\end{pmatrix}.
\]
Subtracting $(a+d)A$ and adding $(ad-bc)I$ cancels each of the four
entries, using only commutativity and distributivity.
For the numerical matrix $H$ above this gives
\[
 H^2=5H+2I,\qquad H^3=27H+10I,\qquad H^4=145H+54I.
\]
Each successive equality is obtained by multiplying the previous one
by $H$ and replacing $H^2$ by $5H+2I$. Over $\mathbb Q$,
\[
 H^{-1}=\tfrac12(H-5I),\qquad
 H^{-2}=\tfrac14(-5H+27I).
\]
The first identity follows from $H(H-5I)=2I$; squaring it and reducing
$H^2$ proves the second. Thus all of the article's elementary numerical
examples have exact integer or rational proofs.


\section{Coefficients, traces, and division-free identities}
\label{sec:trace}

Throughout this section $R$ is a commutative ring with identity and $A$ is an
$n\times n$ matrix.  All sums and coefficient operations below are finite.
Write
\[
 E(t)=\det(I+tA)=\sum_{k=0}^n e_k t^k,
 \qquad s_j=\operatorname{tr}(A^j)\quad(j\geq1),
 \qquad d_k=(-1)^k e_k.
\]
Set $e_0=d_0=1$, and set $e_k=d_k=0$ for $k>n$.  Then
\[
 p_A(X)=X^n+d_1X^{n-1}+\cdots+d_n.
\]
The coefficient $d_k$ here is the coefficient called $c_{n-k}$ in the
article.  A natural-number multiple such as $k e_k$ means repeated addition;
it does not presuppose that $k$ is invertible.

\subsection{Principal minors and the first coefficients}

\begin{proposition}[Principal-minor formula]
\label{prop:trace-principal}
For $0\leq k\leq n$,
\[
 e_k=\sum_{\substack{S\subseteq\{1,\ldots,n\}\\ |S|=k}}
       \det A[S,S].
\]
In particular $e_1=\operatorname{tr}(A)$, $e_n=\det A$, and
\[
 e_2=\sum_{1\leq i<j\leq n}
       (a_{ii}a_{jj}-a_{ij}a_{ji}).
\]
\end{proposition}
\begin{proof}
Expand the determinant by its permutation definition, and in each factor
$\delta_{i,\sigma(i)}+t a_{i,\sigma(i)}$ choose either summand.  A choice of
the second summand on precisely the rows in $S$ contributes a term of degree
$|S|$.  The other factors vanish unless $\sigma$ fixes every index outside
$S$.  Such permutations are exactly the extensions, by fixed points, of
permutations of the increasingly ordered set $S$.  Their signs are the
signs of those restricted permutations: adjoining a fixed point does not
change the parity of a permutation.  Thus their sum is the determinant of
the principal submatrix indexed by $S$.  This gives the asserted coefficient
by a finite reindexing.  The special cases are its instances for subsets of
size one, two, and $n$.
\end{proof}

\subsection{Newton identities without division or eigenvalues}

\begin{proposition}[Finite Jacobi identity]
\label{prop:trace-jacobi}
For a matrix $M(t)$ of polynomials,
\[
 \frac{d}{dt}\det M(t)
 =\operatorname{tr}\bigl(\operatorname{adj} M(t)\,M'(t)\bigr).
\]
The derivative is the formal derivative of a finite coefficient list.
\end{proposition}
\begin{proof}
Apply the product rule to every product in the permutation formula for the
determinant.  Group terms according to the position $(i,j)$ of the
differentiated factor.  The sum of the remaining signed products is the
cofactor $C_{ij}(M)$, so the result is
$\sum_{i,j}C_{ij}(M)M'_{ij}$.  Since
$(\operatorname{adj}M)_{ji}=C_{ij}(M)$, this sum is the stated trace.
Both the product rule and this regrouping are identities between finite
sums, valid in every characteristic.
\end{proof}

\begin{theorem}[Division-free Newton identities]
\label{thm:trace-newton}
For every $k\geq1$,
\begin{equation}
 k e_k=\sum_{i=1}^{k}(-1)^{i-1}e_{k-i}s_i,
 \qquad
 k d_k+\sum_{i=1}^{k}d_{k-i}s_i=0.
 \label{eq:trace-newton}
\end{equation}
These formulas include the cases $k>n$, using the zero convention above.
\end{theorem}
\begin{proof}
Fix $k$.  The finite geometric sum
\[
 U(t)=\sum_{j=0}^{k-1}(-t)^jA^j
\]
satisfies $(I+tA)U(t)=I+(-1)^{k-1}t^kA^k$.  Consequently, modulo $t^k$,
the adjugate identity gives
\[
 \operatorname{adj}(I+tA)=E(t)U(t).
\]
For completeness, multiply
$\operatorname{adj}(I+tA)(I+tA)=E(t)I$ on the right by $U(t)$ and discard
only coefficients of degree at least $k$; this gives the displayed
congruence without cancelling $E(t)$.  The finite Jacobi identity now yields
\[
 E'(t)\equiv E(t)\sum_{j=0}^{k-1}(-1)^j s_{j+1}t^j\pmod{t^k}.
\]
Comparing the coefficient of $t^{k-1}$ proves the first equality.  Multiplying
it by $(-1)^k$ proves the second.  Differentiation was performed on ordinary
polynomials before truncation; no derivative on the quotient ring
$R[t]/(t^k)$ has been assumed.
\end{proof}

This proof does not require a field, characteristic zero, an eigenvalue, an
algebraic closure, or an inverse of $A$.  It also proves the useful adjugate
trace identities
\begin{equation}
 p_A'(X)=\operatorname{tr}\operatorname{adj}(XI-A),\qquad
 Xp_A'(X)-n p_A(X)
   =\operatorname{tr}\bigl(A\operatorname{adj}(XI-A)\bigr).
 \label{eq:trace-adjugate}
\end{equation}
Indeed the first is Jacobi's formula with derivative matrix $I$; the second
is the trace of
$(XI-A)\operatorname{adj}(XI-A)=p_A(X)I$.

\begin{proposition}[A supplied splitting is enough]
\label{prop:trace-splitting}
Suppose a list $\lambda_1,\ldots,\lambda_n$ in a commutative extension ring
is supplied with the identity
$p_A(X)=\prod_{j=1}^n(X-\lambda_j)$.  Then $e_k$ is the $k$th elementary
symmetric expression in that list and, for every $r\geq1$,
\[
 \operatorname{tr}(A^r)=\sum_{j=1}^n\lambda_j^r.
\]
No diagonalization hypothesis is necessary.
\end{proposition}
\begin{proof}
Coefficient comparison gives the assertion about $e_k$.  Put
$u_r=\sum_j\lambda_j^r$.  Differentiate the finite product
$E(t)=\prod_j(1+t\lambda_j)$ and multiply each factor formally by its
truncated geometric inverse.  The same coefficient argument as above gives
$k d_k+\sum_{i=1}^k d_{k-i}u_i=0$.  In this recurrence the coefficient of
$u_k$ is $d_0=1$.  Subtract the recurrence for $s_k$.  Induction on $k$,
using equality of all preceding $u_i$ and $s_i$, gives $u_k=s_k$.
This last induction needs no division by $k$.
\end{proof}

\subsection{Bell-polynomial formulas with all denominators accounted for}

Define the complete exponential Bell polynomials by the finite recurrence
\begin{equation}
 B_0=1,\qquad
 B_k(x_1,\ldots,x_k)=
  \sum_{i=1}^{k}\binom{k-1}{i-1}x_iB_{k-i}(x_1,\ldots,x_{k-i}).
 \label{eq:trace-bell-rec}
\end{equation}
This recurrence is a definition over $\mathbb Z$ and involves no fractions.
For example, $B_1=x_1$, $B_2=x_1^2+x_2$, and
$B_3=x_1^3+3x_1x_2+x_3$.

\begin{theorem}[Integral Bell identity]
\label{thm:trace-bell}
For every $k\geq0$ in every commutative ring,
\begin{equation}
 k!e_k=B_k\bigl(s_1,-s_2,2!s_3,\ldots,
                    (-1)^{k-1}(k-1)!s_k\bigr).
 \label{eq:trace-bell}
\end{equation}
Thus, if $k!$ is a unit, $d_k$ is $(-1)^k/k!$ times the displayed Bell
polynomial.  In particular the article's determinant formula is obtained
at $k=n$ when $n!$ is a unit.
\end{theorem}
\begin{proof}
Use induction, retaining all previously established identities.  Substitute
$x_i=(-1)^{i-1}(i-1)!s_i$ into the Bell recurrence and substitute
$B_{k-i}=(k-i)!e_{k-i}$ from the earlier stages.  The elementary integer
identity
\[
 \binom{k-1}{i-1}(i-1)!(k-i)!=(k-1)!
\]
reduces the result to
$(k-1)!\sum_{i=1}^k(-1)^{i-1}s_i e_{k-i}$, which equals $k!e_k$ by
\eqref{eq:trace-newton}.  There is no cancellation of $(k-1)!$, so the proof
works in positive characteristic and in rings with additive torsion.
\end{proof}

The first four identities are
\begin{align}
 e_1&=s_1,\notag\\
 2e_2&=s_1^2-s_2,\notag\\
 6e_3&=s_1^3-3s_1s_2+2s_3,\label{eq:trace-lowdim}\\
 24e_4&=s_1^4-6s_1^2s_2+3s_2^2+8s_1s_3-6s_4.\notag
\end{align}
They follow by explicitly evaluating the first four Bell recurrences, so
the displayed formulas supply, respectively, the determinant in dimensions
$1,2,3,4$ after the indicated exact division.  They also give the first four
nonleading characteristic coefficients in all larger dimensions.  For
example, if $2$ is invertible, the dimension-three identity is
\[
 A^3-s_1A^2+\frac{s_1^2-s_2}{2}A-\det(A)I=0.
\]
If $2$ and $3$ are invertible, its dimension-four counterpart is
\[
 A^4-s_1A^3+\frac{s_1^2-s_2}{2}A^2
 -\frac{s_1^3-3s_1s_2+2s_3}{6}A+\det(A)I=0.
\]
Over an arbitrary ring, retain the coefficients $e_k$ from principal minors;
one may multiply the displayed equations by the denominators to obtain valid
consequences, but those multiplied equations need not imply Cayley--Hamilton.

\paragraph{Why the characteristic restriction matters.}
For a prime $p$, the two $p\times p$ matrices $0$ and $I$ over
$\mathbb F_p$ have equal traces of every positive power, namely zero, while
their determinants are respectively $0$ and $1$.  Their characteristic
polynomials are $X^p$ and $X^p-1$.  Thus the traces alone cannot recover the
characteristic coefficients in arbitrary characteristic.  The factor
$k!$ in \eqref{eq:trace-bell} cannot simply be cancelled.

\subsection{The Hessenberg determinant formula}

Let $H_k$ be the $k\times k$ matrix with entries, in one-based indexing,
\[
 (H_k)_{ij}=
 \begin{cases}
 s_{i-j+1},&j\leq i,\\
 k-i,&j=i+1,\\
 0,&j>i+1.
 \end{cases}
\]
Thus its first row is $(s_1,k-1,0,\ldots,0)$ and its last row is
$(s_k,s_{k-1},\ldots,s_1)$.  Put $D_0=1$ and $D_k=\det H_k$.

\begin{theorem}[Integral Hessenberg identity]
\label{thm:trace-hessenberg}
For every $k\geq0$, $D_k=k!e_k$.  In particular, when $k!$ is a unit,
\[
 c_{n-k}=d_k=\frac{(-1)^k}{k!}\det H_k\qquad(0\leq k\leq n).
\]
\end{theorem}
\begin{proof}
Expand $D_k$ along its first column.  In the minor corresponding to row $i$,
the first $i-1$ rows force the entries
$k-1,k-2,\ldots,k-i+1$ successively: the first such row has only its first
available entry nonzero, and the same holds after each removal.  The
remaining lower-right matrix is $H_{k-i}$.  Therefore
\[
 D_k=\sum_{i=1}^k(-1)^{i-1}s_i
       \left(\prod_{h=1}^{i-1}(k-h)\right)D_{k-i}.
\]
The empty product is $1$.  Inductively substituting
$D_{k-i}=(k-i)!e_{k-i}$ turns the product of the two factorial factors into
$(k-1)!$.  Newton's identity now gives $D_k=k!e_k$, as required.
\end{proof}

\subsection{A correctly initialized Faddeev--LeVerrier algorithm}

Assume for now that $1,2,\ldots,n$ are units of $R$.  Define
\begin{equation}
 F_0=I,\qquad
 d_k=-k^{-1}\operatorname{tr}(AF_{k-1}),\qquad
 F_k=AF_{k-1}+d_k I\quad(1\leq k\leq n).
 \label{eq:trace-flv}
\end{equation}
These assignments compute the characteristic coefficients.  Indeed the
invariant
\[
 F_k=A^k+d_1A^{k-1}+\cdots+d_kI
\]
is immediate by induction.  Assuming the coefficients through $k-1$ have
been computed correctly, Newton's identity gives
\[
 \operatorname{tr}(AF_{k-1})
  =s_k+d_1s_{k-1}+\cdots+d_{k-1}s_1=-k d_k,
\]
which proves the next assignment.  The invariant at $k=n$ gives
$F_n=p_A(A)=0$.  The adjugate coefficient formula additionally gives
$\operatorname{adj}(A)=(-1)^{n-1}F_{n-1}$ for $n\geq1$; if $d_n$ is a unit,
$A^{-1}=-d_n^{-1}F_{n-1}$.

For integer matrices the same algorithm uses exact integer division by
the positive number $k$.  Its divisibility is not guessed: the principal-minor
definition supplies the integer witness $-d_k$ for
$\operatorname{tr}(AF_{k-1})=k(-d_k)$.  Integer quotient and remainder are
primitive recursive, and their defining equations verify each step.  More
generally an effective exact-division operation with this correctness
property suffices; mere characteristic zero does not supply such an operation.

The article's displayed recurrence with initial value $M_0=0$ contains
$1/(k-1)$ at $k=1$.  Read literally it cannot be initialized.  One must give
$M_1=I$ and start that recurrence at $k=2$, or use
\eqref{eq:trace-flv} from $F_0=I$.

\subsection{Finite formal logarithms, exponentials, and the partition formula}

Fix a precision $N$ and suppose $1,2,\ldots,N$ are units, as they are in a
$\mathbb Q$-algebra.  Every expression in this subsection is reduced modulo
$z^{N+1}$.  Define
\[
 L_N(z)=-\sum_{j=1}^N\frac{s_j}{j}z^j,
 \qquad
 G_N(z)=\sum_{r=0}^N\frac{L_N(z)^r}{r!}\pmod{z^{N+1}}.
\]
These are finite polynomial definitions, not limits.  Then
\begin{equation}
 \det(I-zA)\equiv G_N(z)
   =\exp_N\bigl(\operatorname{tr}\log_N(I-zA)\bigr)
       \pmod{z^{N+1}}.
 \label{eq:trace-explog}
\end{equation}
To prove the identity, write $g_k$ for the coefficient of $z^k$ in $G_N$.
Differentiating the ordinary finite sum defining $G_N$ gives
$G_N'\equiv L_N'G_N\pmod{z^N}$: the omitted term is
$L_N'L_N^N/N!$ and has degree at least $N$.  Coefficient comparison therefore
gives
\[
 g_0=1,\qquad k g_k=-\sum_{j=1}^k s_j g_{k-j}\quad(1\leq k\leq N).
\]
The coefficients $d_k$ of $\det(I-zA)$ satisfy exactly this recurrence by
Newton's identity.  Induction and multiplication by the supplied inverse of
$k$ give $g_k=d_k$ through degree $N$.  This proves
\eqref{eq:trace-explog} without analytic convergence or a logarithm branch.
No nonzero-eigenvalue assumption is needed.  If $N>n$, the coefficients of
degrees $n+1,\ldots,N$ in $G_N$ consequently vanish.

Expanding the finite exponential and grouping its products according to
the number $m_j$ of occurrences of $s_jz^j/j$ gives the useful corrected
partition formula
\begin{equation}
 d_k=
 \sum_{\substack{m_1,\ldots,m_k\geq0\\
                 \sum_{j=1}^k j m_j=k}}
 \prod_{j=1}^k
   \frac{(-1)^{m_j}s_j^{m_j}}{j^{m_j}m_j!}.
 \label{eq:trace-partition}
\end{equation}
For a fixed tuple put $r=\sum_jm_j$.  Its terms in $L_N^r$ occur
$r!/\prod_jm_j!$ times; division by $r!$ proves the displayed coefficient.
Every index has the bound $m_j\leq k$, so this is a finite bounded sum.
The article's product with sign $(-1)^{m_j+1}$ for each of $n$ factors differs
from the correct sign by $(-1)^n$; for odd $n$ it would already give the
incorrect leading coefficient $c_n=-1$.  Equation
\eqref{eq:trace-partition} uses the corrected sign.  In rings where the
fractions are unavailable, the integral Bell identity is the applicable
statement instead.

The familiar expression
$p_A(\lambda)=\lambda^n\exp(\operatorname{tr}\log(I-A/\lambda))$
can therefore be read rigorously, at every finite precision, by putting
$z=\lambda^{-1}$ in \eqref{eq:trace-explog}.  This is a formal Laurent
calculation.  It does not assert the existence of a numerical matrix
logarithm for every nonzero value of $\lambda$.

\subsection{The exact induction obligations for iPRA}

For the arithmetic implementation, subsets are bit masks below $2^n$,
permutations are bounded lists, and polynomials are finite coefficient
arrays with an explicit length.  Traces, determinants, products, powers,
coefficient selection, formal derivatives, and truncation are primitive
recursive operations on these codes.  The earlier ring-identity
interpretation is used for the universal ring assertions.  Rational
instances use signed integer numerator--positive denominator pairs;
statements about supplied inverses include their finite verification
equations.

Computability by itself is not the proof of iPRA-provability.  The preceding
proofs use the following explicit induction invariants, each a decidable
equality of finite arrays or a bounded conjunction of such equalities.
For a coded ring presentation use its decidable equality; for generic
identities compare the corresponding normalized integer polynomial arrays.
\begin{enumerate}
\item The geometric-sum invariant is
$(I+tA)\sum_{j=0}^{q}(-t)^jA^j
 =I+(-1)^qt^{q+1}A^{q+1}$.
Its successor step is obtained by adjoining one term and cancelling the
two adjacent terms of degree $q+1$.
\item Determinant differentiation is a bounded sum of product-rule
identities, and the principal-minor argument is the explicit finite
bijection $(\sigma,S)\leftrightarrow(S,\sigma|_S)$ on nonzero summands.
Finite-sum reindexing is proved by induction on the enumerated list length;
zero summands are retained, so no existence search for a nonzero entry occurs.
\item In the Bell proof retain the table invariant
$B_j=j!e_j$ for every $j\leq q$.  Its next row follows by
\eqref{eq:trace-bell-rec} and \eqref{eq:trace-newton}.  For the Hessenberg
proof retain $D_j=j!e_j$ on the same prefix.  Their factorial identities are
equalities of finite products and binomial coefficients, with the usual
Pascal recursion giving quantifier-free induction proofs.
\item In the Faddeev--LeVerrier proof retain the polynomial expression for
$F_q$ and correctness of the first $q$ coefficients.  The next numerator
is exactly $-k d_k$ by Newton's identity, so unit multiplication or verified
exact integer division proves the next step.
\item In the finite exponential proof retain $g_j=d_j$ for $j\leq q$.
The displayed coefficient recurrence and the explicit inverse of $q+1$
give the successor step.  The partition expansion is a bounded grouping of
words of length at most $N$ by their occurrence-count vectors.
\end{enumerate}
A prefix invariant can be represented by one primitive recursive Boolean
conjunction over its bounded list.  Induction on that Boolean equation is
quantifier-free induction, so no induction over an unbounded quantified
formula is being silently introduced.  These invariants give the concrete
proof obligations for the object-language formalization, rather than
inferring provability merely from the existence of the algorithms.

% Sections prepared independently for inclusion in the main manuscript.
% Requires amsmath, amssymb, and theorem/proposition/lemma/corollary/remark environments.
% No custom macros are used.

\section{Finite consequences: powers and minimal polynomials}
\label{sec:finite-applications}

Throughout this section, $n\geq1$ and
\[
 p(X)=X^n+\sum_{i=0}^{n-1}c_iX^i=\det(XI-A).
\]
All ring statements below are equational statements, or implications from
explicit finite data. To read them internally in intuitionistic primitive
recursive arithmetic, fix a primitive recursive presentation of the ring
whose ring laws have iPRA proofs. The field algorithms additionally require
primitive recursive equality and inversion of nonzero elements, with
iPRA proofs of the field axioms and the guarded inverse identities for
that fixed presentation. Rational
numbers and finite fields provide concrete examples. This convention does
not quantify over arbitrary rings as objects of first-order arithmetic.

\subsection{Monic division without a zero test}

\begin{lemma}[Fixed-length monic division]
For every coefficient list $f$ and monic $p$ of displayed degree $n$, there
are primitive recursively computed lists $q,r$ such that
\[
 f=qp+r,\qquad r_i=0\quad(i\geq n).
\]
They are unique as polynomials. No division by a ring element is used.
\end{lemma}
\begin{proof}
Use the input list length as an upper bound for degree. Starting at its last
index $d\geq n$, subtract $f_dX^{d-n}p$ and add $f_dX^{d-n}$ to the quotient.
The coefficient at $d$ becomes zero, regardless of whether it was already
zero. Continue down to $n$. The invariant is the coefficientwise identity
$f=q_{\mathrm{current}}p+r_{\mathrm{current}}$ with all processed high
coefficients zero. It is an equality of bounded lists and is preserved by
one distributive-law calculation.

For uniqueness, suppose $p(q-q')=r'-r$. Look at coefficients from the fixed
maximum list length downwards. The highest not-yet-eliminated coefficient
of the left side is the corresponding coefficient of $q-q'$, since the
leading coefficient of $p$ is $1$; the right side is zero in every degree
at least $n$. Descending induction therefore sets every coefficient of
$q-q'$ to zero, and then every coefficient of $r-r'$ to zero. This argument
needs neither an integral domain nor a decision of whether a coefficient
is zero.
\end{proof}

Evaluation in a matrix is multiplicative for scalar polynomials because
scalar matrices commute with $A$: expansion of the finite double sum gives
$(fg)(A)=f(A)g(A)$. Consequently $p(A)=0$ implies $f(A)=r(A)$.
For powers, one can avoid constructing the quotient. Let $r_{0}=1$ and
write $r_k=\sum_{i<n}r_{k,i}X^i$. Then
\begin{align*}
 r_{k+1,0}&=-c_0r_{k,n-1},\\
 r_{k+1,i}&=r_{k,i-1}-c_ir_{k,n-1}\qquad(1\leq i<n).
\end{align*}
The induction invariant is $A^k=\sum_{i<n}r_{k,i}A^i$.
Thus every nonnegative power has a computed representative of degree
less than $n$. Reduced multiplication followed by binary exponentiation
also computes this representative; binary exponentiation improves
efficiency but is unnecessary for primitive recursiveness.

Negative powers are also covered: use the explicit inverse from
Theorem~\ref{thm:inverse}, take its positive powers, and apply the same
monic reduction. The exact numerical examples were proved in
Section~\ref{sec:main}.

\subsection{Computing a minimal polynomial by bounded elimination}

\begin{proposition}[Effective field version]
Over a nontrivial discrete field with primitive recursive operations and
inversion, a primitive recursive algorithm constructs a monic polynomial
$m_A$ satisfying
\[
 m_A(A)=0,\qquad \deg m_A\leq n,\qquad
 h(A)=0\ \Longleftrightarrow\ m_A\mid h.
\]
In particular $m_A\mid p$; conversely that divisibility implies $p(A)=0$.
\end{proposition}
\begin{proof}
Flatten $I,A,A^2,\ldots$ into vectors $v_j$ of length $n^2$.
For $d=1,\ldots,n^2$, solve
\[
 [v_0\ \cdots\ v_{d-1}]x=-v_d
\]
by row reduction. Search for a pivot only within the remaining finite
rectangle. A nonzero pivot is inverted and its row is used to eliminate
its column. Every operation is invertible, so the transformed system is
equivalent to the original one. If a zero row of coefficients has nonzero
right-hand side, the stored product of row operations supplies a row
witness $w$ with $w[v_0\cdots v_{d-1}]=0$ but $wv_d\ne0$. Otherwise set
the free variables to zero and recover a solution by back substitution.
There is no search over possible field coefficients.

Some $d\leq n^2$ succeeds: row reduction of the $n^2\times(n^2+1)$
matrix $[v_0\cdots v_{n^2}]$ has at most $n^2$ pivots, so a free column
produces a dependence. Normalize its last nonzero coefficient. Since
$I\ne0$, that last index is positive. Select the first successful $d$
and its coefficients, obtaining $m_A=X^d+\sum_{j<d}x_jX^j$.
The supplied inconsistency witnesses for earlier $d$ certify minimality.
Cayley--Hamilton makes the system successful already at $d=n$, so the
output has $d\leq n$.

Divide any annihilator $h$ by $m_A$: $h=qm_A+r$, with $\deg r<d$.
Then $r(A)=0$. A bounded scan either verifies every coefficient of $r$
is zero, or finds its last nonzero coefficient. In the latter case,
division by that coefficient makes a monic annihilator of degree less
than $d$, contradicting the stored minimality certificates. A nonzero constant
remainder would imply $I=0$, which is excluded by the nontrivial field
and $n\ge1$ hypotheses. Thus $r=0$.
The converse follows from multiplicativity of evaluation. This also
proves uniqueness of the monic generator.
\end{proof}

The loops are bounded by displayed dimensions; their invariants are
equalities and decidable inequalities between finite arrays. A field
presented only by an abstract classical existence axiom for inverses
does not supply the algorithm assumed here. Nor is a minimal-polynomial
algorithm over an arbitrary commutative ring asserted.

\subsection{Companion matrices}

For a monic $f=X^n+a_{n-1}X^{n-1}+\cdots+a_0$, define its companion matrix
by
\[
 C_fe_i=e_{i+1}\ (i<n-1),\qquad
 C_fe_{n-1}=-\sum_{i<n}a_ie_i.
\]
Then $C_f^ie_0=e_i$ for $i<n$, so $f(C_f)e_0=0$. Commuting with
$C_f^i$ shows $f(C_f)e_i=0$ for every basis vector, and hence $f(C_f)=0$.
Expansion of $\det(XI-C_f)$ down the last column gives $f$: the cofactor
at row $i$ of that column is $X^i$. Indeed, its minor has $i$ diagonal
$X$ factors and $n-1-i$ subdiagonal $-1$ factors; its cofactor sign
cancels the latter signs. The final row contributes
$(X+a_{n-1})X^{n-1}$.

More strongly, over every commutative ring,
\[
 h(C_f)=0\quad\Longleftrightarrow\quad f\mid h.
\]
Divide $h=qf+r$ with $r$ supported below $n$. If $h(C_f)=0$, then
$0=r(C_f)e_0=\sum_{i<n}r_ie_i$, which says every $r_i$ is zero.
This argument does not require a field or linear dependence search.

\section{Eigenvalue statements with explicit finite data}

If $Av=\lambda v$, induction gives $A^kv=\lambda^kv$, and therefore
\[
 f(A)v=f(\lambda)v
\]
for every scalar polynomial $f$. With $f=p$, this yields $p(\lambda)v=0$.
Over a discrete field, a nonzero vector has a nonzero component found by
a bounded scan; inversion of that component gives $p(\lambda)=0$.
Over a ring the same conclusion holds if a component is supplied with a
unit inverse. Without such a hypothesis it can fail: in $\mathbb Z/6$,
$A=(2)$, $\lambda=5$, and $v=(2)$ satisfy $Av=\lambda v$ although
$p(5)=3\ne0$.

Conversely, over the effective field, $p(\lambda)=0$ means that
$\lambda I-A$ has zero determinant. Row reduction cannot have a pivot
in every column, because elementary row operations change determinant
by known nonzero factors and a triangular matrix with all pivots has
nonzero determinant. Set one free variable to $1$, the others to $0$,
and back substitute. The result is a certified nonzero eigenvector.
This proves the eigenvalue/root equivalence constructively under the
stated field assumptions.

Suppose matrices $S,T$ with $ST=TS=I$ are supplied and $B=TAS$.
Finite multiplication gives $B^k=TA^kS$, while determinant
multiplicativity gives
\[
 \det(XI-B)=\det T\,\det(XI-A)\,\det S=p(X).
\]
If $B$ is upper triangular with diagonal $\lambda_0,\ldots,\lambda_{n-1}$,
its determinant polynomial is $\prod_i(X-\lambda_i)$. To verify this
directly, a nonzero Leibniz summand would require $\sigma(i)\geq i$ for
every $i$; equality of the two finite sums of indices forces all
$\sigma(i)=i$. Products of triangular matrices multiply their diagonal
entries, so $\operatorname{tr}(B^k)=\sum_i\lambda_i^k$.
Finally $\operatorname{tr}(UV)=\operatorname{tr}(VU)$ follows by
interchanging two finite sums, and hence
\[
 \operatorname{tr}(A^k)=\sum_i\lambda_i^k.
\]
These are proofs from a supplied triangularization certificate; they do
not assume an algorithm finding roots or a Jordan form. A supplied
diagonalization is the special case in which $B$ is diagonal.

\section{Finitely generated modules and integral elements}

\subsection{The determinant trick without a basis}

\begin{proposition}[Finite generator form]
Let $g_0,\ldots,g_{n-1}$ be elements of an $R$-module, and let $\phi$ be
$R$-linear. Suppose
\[
 \phi(g_j)=\sum_{i<n}a_{ij}g_i.
\]
Writing $p=\det(XI-A)$, one has $p(\phi)(g_j)=0$ for every $j$, and
therefore $p(\phi)(\sum_jx_jg_j)=0$ for every supplied coordinate list
$x$. If these elements generate the module, $p(\phi)=0$ on it.
\end{proposition}
\begin{proof}
Let $\pi(x)=\sum_jx_jg_j$. The supplied equations say
$\phi\pi=\pi A$. Induction gives $\phi^k\pi=\pi A^k$; addition with the
coefficients of $p$ gives $p(\phi)\pi=\pi p(A)=0$. Apply to the unit
vectors, or to any supplied $x$. Linear independence of the generators
has never been used.
\end{proof}

In arithmetic the last assertion is read using a coordinate witness for
the element in question, or a primitive recursive coordinate function
when a uniform computational presentation is desired. It is not an
invocation of a choice principle to select coordinates in an arbitrary
abstract module.

If every $a_{ij}$ belongs to an ideal $J$, then the coefficient of
$X^{n-k}$ in $p$ belongs to $J^k$: the principal-minor formula displays
it as a finite sum of products of $k$ such entries. This proves the
usual strengthened determinant trick with explicit membership
certificates.

For a certificate version of Nakayama's lemma, suppose
$g_j=\sum_i a_{ij}g_i$ with $a_{ij}\in J$. Apply the proposition with
$\phi=1$. The scalar
$\delta=\det(I-A)=1-b$, for a displayed $b\in J$, annihilates every
$g_j$. If an inverse $u\delta=1$ is supplied, multiplying gives $g_j=0$.
Thus the whole generated module is zero. The familiar assumption that
$J$ lies in the Jacobson radical supplies this unit property in the
abstract formulation; the finite arithmetic theorem uses the inverse
as data and does not search for it.

\subsection{Closure of integral elements with a degree bound}

\begin{proposition}[Explicit integrality construction]
In a commutative $R$-algebra, suppose $\alpha_j$ satisfies a supplied
monic polynomial $f_j$ of positive degree $d_j$, for $1\leq j\leq k$.
For every supplied polynomial expression $\beta$ in the $\alpha_j$,
one can construct a monic annihilator of $\beta$ of degree
$D=\prod_{j=1}^k d_j$.
\end{proposition}
\begin{proof}
Use the $D$ monomials
$g_e=\prod_j\alpha_j^{e_j}$, $0\leq e_j<d_j$, including $g_0=1$.
Reduce high powers by the supplied monic relations, one variable at a
time. Every reduction preserves the represented element and decreases
the exponent currently being processed. It computes coefficients for
$\beta g_e$ in the displayed list. Let $A$ be the resulting $D\times D$
matrix. Multiplication by $\beta$ satisfies the generator equations,
so the determinant trick gives $p_A(\beta)g_0=0$. As $g_0=1$,
$p_A(\beta)=0$. The construction uses lists and bounded monic division.
It neither assumes the monomials are independent nor computes a basis
of their span.
\end{proof}

Taking $R=\mathbb Z$ proves closure of integral elements under addition,
multiplication, and negation, with a computed monic equation. In
particular two input equations of degrees $d,e$ give an equation of
degree $de$ for both the sum and the product. This bypasses any theorem
that a submodule of a finitely generated module is finitely generated.

For an element $\alpha$ of a finite extension presented by a
$\mathbb Q$-basis, a unit coordinate vector, and a rational multiplication
table, form the matrix $L_\alpha$ of multiplication by $\alpha$.
Associativity and the unit give
$h(L_\alpha)=L_{h(\alpha)}$ and
$L_{h(\alpha)}(1)=h(\alpha)$. Thus its matrix minimal polynomial is
exactly the minimal polynomial of $\alpha$. The bounded algorithm above
computes it. Computing the characteristic polynomial alone only gives
an annihilator; it need not be minimal. For example the multiplication
matrix of $1$ in any extension of degree $D>1$ has characteristic
polynomial $(X-1)^D$ and minimal polynomial $X-1$.

If $\alpha$ is integral over $\mathbb Z$, its rational minimal polynomial
has integer coefficients. Here is the finite algebra needed for that
extra claim. Let $P\in\mathbb Z[X]$ be a monic annihilator and
$P=mq$ in $\mathbb Q[X]$, with $m,q$ monic. Clear denominators and
divide coefficient gcds to write $f=am$, $g=bq$ as primitive integer
polynomials with positive leading coefficients $a,b$. These are positive
integers: the coefficient gcd divides the leading coefficient of the
cleared monic polynomial. A product of
primitive integer polynomials is primitive: otherwise a prime dividing
all its coefficients, found by bounded search from their gcd, would
make its reduction zero modulo that prime. The two factors each have
a last nonzero reduced coefficient; their product gives a nonzero
coefficient in the sum of these two degrees, a contradiction. Hence
$fg=abP$ has content $1$, whereas the right side has content $ab$.
Therefore $a=b=1$ and $m\in\mathbb Z[X]$. Integer gcd, bounded prime
search, and finite-field multiplication all have primitive recursive
implementations.

\section{Interpolation from supplied splitting data}

Suppose a monic annihilating polynomial is supplied in the form
\[
 p(X)=\prod_{j<r}(X-\lambda_j)^{m_j},\qquad
 m_j\geq1,\quad\sum_jm_j=n,
\]
and every difference $\lambda_i-\lambda_j$ for $i\ne j$ has a supplied
unit inverse. These assumptions are stronger than the mere existence of
roots and make all interpolation operations explicit.

\subsection{Distinct roots and Sylvester's formula}

If all $m_j=1$, set
\[
 L_j(X)=\prod_{i\ne j}\frac{X-\lambda_i}{\lambda_j-\lambda_i}.
\]
For every polynomial $f$,
\[
 \operatorname{rem}_p f=\sum_j f(\lambda_j)L_j(X),\qquad
 f(A)=\sum_jf(\lambda_j)L_j(A).
\]
To prove uniqueness, a polynomial vanishing at $\lambda_0$ is divisible
by $X-\lambda_0$, by monic division. If it also vanishes at $\lambda_1$,
evaluate the quotient identity there and cancel the supplied unit
$\lambda_1-\lambda_0$ to obtain the next factor. Continuing shows that
a polynomial of degree less than $n$ vanishing at every root is zero.
The displayed candidate has the required values, so it is the
remainder. Evaluating its congruence modulo $p$ gives the matrix formula.
The same argument applied to products shows
$L_i(A)L_j(A)=0$ for $i\ne j$, $L_j(A)^2=L_j(A)$, and
$\sum_jL_j(A)=I$.

\subsection{Repeated roots: jets and Hasse derivatives}

The $s$th Hasse derivative at $\lambda$ is the coefficient of $Z^s$ in
$f(\lambda+Z)$:
\[
 D^{[s]}f(\lambda)=\sum_{k\geq s}\binom{k}{s}f_k\lambda^{k-s}.
\]
This is a finite formula in every characteristic. Ordinary derivatives
satisfy $f^{(s)}(\lambda)=s!D^{[s]}f(\lambda)$, so ordinary derivative
conditions are equivalent only when the relevant factorials are units.

Let $q_j=p/(X-\lambda_j)^{m_j}$. Write
$q_j(\lambda_j+Z)=\sum_s a_{j,s}Z^s$ and put
\[
 b_{j,0}=a_{j,0}^{-1},\qquad
 b_{j,t}=-a_{j,0}^{-1}\sum_{s=1}^t a_{j,s}b_{j,t-s}
 \quad(1\leq t<m_j).
\]
The constant $a_{j,0}=\prod_{i\ne j}(\lambda_j-\lambda_i)^{m_i}$
has a supplied inverse. Convolution proves
$q_j(\lambda_j+Z)\sum_{t<m_j}b_{j,t}Z^t\equiv1\pmod{Z^{m_j}}$.
Thus
\[
 E_j(X)=q_j(X)\sum_{t<m_j}b_{j,t}(X-\lambda_j)^t
\]
is $1$ modulo $(X-\lambda_j)^{m_j}$ and $0$ modulo every other factor.
For arbitrary jet data $d_{j,s}$ the polynomial
\[
 r=\operatorname{rem}_p\left(
 \sum_jE_j(X)\sum_{s<m_j}d_{j,s}(X-\lambda_j)^s\right)
\]
therefore has exactly these jets and degree less than $n$.

For uniqueness, note that the same truncated inverse construction shows
that any product of the other factors is invertible modulo
$(X-\lambda_j)^{m_j}$. If a polynomial is divisible by that product
and by $(X-\lambda_j)^{m_j}$, its quotient by the former is divisible by
the latter: multiply its zero residue by this explicit inverse.
Induction therefore makes a polynomial divisible by each factor
divisible by their product $p$. A difference of two candidates of
degree less than $n$ must vanish by monic division uniqueness.
Taking $d_{j,s}=D^{[s]}f(\lambda_j)$ proves $f(A)=r(A)$.

This explains why repeated eigenvalues require derivative data in
characteristic zero, and why repeated values alone do not determine a
unique remainder. The proof also supplies a division-free replacement
for derivatives in positive characteristic, apart from the explicitly
assumed inverse root differences.

\section{Matrix functions as controlled finite calculations}

An arbitrary analytic function is not a first-order arithmetic object.
Here the precise replacement is a specified primitive recursive sequence
of rational coefficients and a primitive recursive modulus for finite
tails. Equalities of the corresponding Cauchy names mean explicit
finite inequalities at every precision. The external real-number
interpretation is useful notation, not an additional completeness axiom
inside iPRA.

\subsection{A sufficient quantitative polynomial-reduction theorem}

Let $A$ be a rational matrix and let $r_k=\operatorname{rem}_pX^k$ as
above. For a polynomial supported below $n$, write
$|r|_1=\sum_{i<n}|r_i|$, and choose rational
\[
 K\geq\max\left(1,\sum_{i<n}|c_i|\right),\qquad
 H\geq\max(1,\|A\|_\infty^{n-1}),
\]
where $\|A\|_\infty=\max_i\sum_j|a_{ij}|$.
The recurrence for $r_k$ gives
$|r_{k+1}|_1\leq K|r_k|_1$, so $|r_k|_1\leq K^k$.
The finite triangle inequality and submultiplicativity of the row norm
give $\|r(A)\|_\infty\leq H|r|_1$.

Suppose a primitive recursive rational sequence $a_k$ and a primitive
recursive function $N$ are supplied together with proofs that
\[
 \sum_{k=N(s)}^m |a_k|K^k\leq2^{-s}
 \qquad(m\geq N(s)).                                      \tag{*}
\]
Define finite approximants
\[
 \beta_{i,M}=\sum_{k<M}a_kr_{k,i},\qquad
 F_M=\sum_{k<M}a_kA^k.
\]
Their exact finite identity is
$F_M=\sum_{i<n}\beta_{i,M}A^i$. Condition (*) bounds every coefficient
tail by $2^{-s}$ and every matrix tail by $H2^{-s}$. It follows, in the
precise sense of these finite estimates, that
\[
 f(A)=\sum_{i<n}\beta_iA^i,\qquad
 \beta_i=\sum_{k\geq0}a_kr_{k,i}.
\]
For a requested matrix error $2^{-e}$, choose an integer $h$ with
$H\leq2^h$ and use $N(e+h)$. Thus the conclusion contains an explicit
primitive recursive modulus. It does not assert that every analytic
function near the spectrum satisfies (*) for this possibly larger $K$.

The statement $f=qp+r$ with both $q$ and $r$ polynomials cannot hold
for a nonpolynomial analytic $f$. Even when an analytic quotient is
available, that is a different assertion. Our argument uses only
polynomial identities for finite partial sums and the specified tail
bounds.

\subsection{An explicit factorial tail bound}

For a positive integer $U$ and $s\geq0$, put
\[
 M(U,s)=2U+2U^2+s+2.
\]
For $k\geq2U$, the ratio of successive terms $U^k/k!$ is at most
$1/2$. Also $U^{2U}/(2U)!\leq U^{2U}\leq2^{2U^2}$, since
$U\leq2^U$. Therefore
\[
 \frac{U^{M(U,s)}}{M(U,s)!}\leq2^{-s-2},\qquad
 \sum_{k=M(U,s)}^v\frac{U^k}{k!}\leq2^{-s-1}
 \quad(v\geq M(U,s)).
\]
The last inequality is the finite geometric-sum bound. All denominators
can be cleared, leaving inequalities of natural numbers proved by
bounded summation and quantifier-free induction.

For rational $t$, choose an integer $U\geq\max(1,|t|K)$.
The exponential coefficients $a_k=t^k/k!$ satisfy (*) with this $M$;
the sine and cosine coefficient subsequences satisfy the same bound.
If a direct matrix-series bound is also wanted, enlarge $U$ to exceed
$|t|\|A\|_\infty$. Consequently matrix exponential, sine, and cosine
have primitive recursive rational Cauchy names and polynomial
representations of degree less than $n$, with the displayed moduli.

\subsection{The article's examples and rotation identities}

For $T=\left(\begin{smallmatrix}1&2\\0&3\end{smallmatrix}\right)$,
induction gives
\[
 T^k=\begin{pmatrix}1&3^k-1\\0&3^k\end{pmatrix}\qquad(k\geq0).
\]
Hence every polynomial obeys
\[
 f(T)=\begin{pmatrix}f(1)&f(3)-f(1)\\0&f(3)\end{pmatrix}.
\]
Apply this exact identity to exponential or sine partial sums and use
the factorial estimate. For every rational $t$, as identities of the
specified Cauchy names,
\[
 e^{tT}=\begin{pmatrix}e^t&e^{3t}-e^t\\0&e^{3t}\end{pmatrix},\qquad
 \sin(tT)=\begin{pmatrix}\sin t&\sin3t-\sin t\\0&\sin3t\end{pmatrix}.
\]
The corresponding coefficients in $c_0I+c_1T$ are
$c_0=(3f(1)-f(3))/2$ and $c_1=(f(3)-f(1))/2$.

More generally the relation $J^2=-I$ gives
$J^{2r}=(-1)^rI$ and $J^{2r+1}=(-1)^rJ$, whence
\[
 e^{tJ}=(\cos t)I+(\sin t)J.
\]
For $J=\left(\begin{smallmatrix}0&1\\-1&0\end{smallmatrix}\right)$
this is the article's two-dimensional rotation matrix. If $P^2=I$,
the same calculation over rational complex numbers gives
$e^{itP}=(\cos t)I+i(\sin t)P$. In particular, with the usual Pauli
matrices
\[
 \sigma_1=\begin{pmatrix}0&1\\1&0\end{pmatrix},\qquad
 \sigma_2=\begin{pmatrix}0&-i\\i&0\end{pmatrix},\qquad
 \sigma_3=\begin{pmatrix}1&0\\0&-1\end{pmatrix},
\]
direct multiplication gives
$\sigma_j^2=I$ and
$\sigma_j\sigma_k+\sigma_k\sigma_j=0$ for $j\ne k$.
Thus $P=\sum_j u_j\sigma_j$ has $P^2=I$ whenever the supplied rational
coordinates satisfy $\sum_j u_j^2=1$.

For rational complex matrices, use the scalar norm
$|a+bi|_1=|a|+|b|$ and its induced row-sum matrix norm. Direct expansion
proves submultiplicativity, so the same factorial estimates justify
the complex-valued formulas above.

For a rational unit vector $u=(u_1,u_2,u_3)$ put
\[
 K=\begin{pmatrix}0&-u_3&u_2\\u_3&0&-u_1\\-u_2&u_1&0\end{pmatrix}.
\]
Entrywise multiplication gives $K^2=uu^{\mathsf T}-I$ and $Ku=0$,
so $K^3=-K$. Hence $K^{2r+1}=(-1)^rK$ and
$K^{2r+2}=(-1)^rK^2$, yielding Rodrigues' formula
\[
 e^{tK}=I+(\sin t)K+(1-\cos t)K^2.
\]
Each displayed analytic identity is obtained from a finite identity
between partial sums followed by the same explicit tail estimates.
Here $\cos t$ and $\sin t$ denote the Cauchy names given respectively by
the finite sums of $(-1)^rt^{2r}/(2r)!$ and
$(-1)^rt^{2r+1}/(2r+1)!$; the exponential uses $t^r/r!$.
No extraction of eigenvalues, analytic continuation, or appeal to a
matrix Lie group is involved. Rational parameters give a fully
arithmetical family; extending it to further real inputs requires
their approximation codes and error bounds as additional data.

\section{Quaternion and octonion boundaries}

\subsection{A precise quaternionic Cayley--Hamilton theorem}

For a rational quaternion $q=a+bi+cj+dk$ define
\[
 L(q)=\begin{pmatrix}
 a&-b&-c&-d\\
 b&a&-d&c\\
 c&d&a&-b\\
 d&-c&b&a
 \end{pmatrix}.
\]
The columns are the coordinates of $q$, $qi$, $qj$, and $qk$.
The quaternion multiplication table proves
$L(qr)=L(q)L(r)$, $L(q+r)=L(q)+L(r)$, and $L(1)=I_4$.
The first column recovers $q$, so $L$ is injective.
Apply $L$ blockwise to $A\in M_n(\mathbb H_{\mathbb Q})$, producing
$L_n(A)\in M_{4n}(\mathbb Q)$. Matrix multiplication is preserved because
each block product is a finite sum of products preserved by $L$.
If $p=\det(XI-L_n(A))$, then
\[
 L_n(p(A))=p(L_n(A))=0.
\]
Injectivity gives $p(A)=0$. This is an explicit monic rational
annihilator of degree $4n$, not a naive quaternionic Leibniz determinant.

There is a degree $2n$ refinement. In $\mathbb Q(i)$ set
\[
 \rho(q)=\begin{pmatrix}a+bi&c+di\\-c+di&a-bi\end{pmatrix},\qquad
 J=\begin{pmatrix}0&1\\-1&0\end{pmatrix}.
\]
Multiplication of the four basis matrices verifies the quaternion
relations, so $\rho$ is an injective algebra homomorphism. Furthermore
$J\overline{\rho(q)}J^{-1}=\rho(q)$. For the block matrix
$B=\rho_n(A)$ and block diagonal $J_n$, this gives
$J_n\overline B J_n^{-1}=B$. Determinant similarity therefore makes
every coefficient of $\det(XI-B)$ fixed by complex conjugation.
A rational complex number fixed by conjugation has zero imaginary
coordinate, so this characteristic polynomial lies in $\mathbb Q[X]$.
Its degree is $2n$ and its evaluation at $A$ is zero by injectivity.
All representation identities involve finitely many rational coordinate
equalities; no quaternionic spectral theorem is required.

For split quaternions with $i^2=j^2=1$ and $ij=k=-ji$, the representation
\[
 a+bi+cj+dk\longmapsto
 \begin{pmatrix}a+b&c+d\\c-d&a-b\end{pmatrix}
\]
is faithful over $\mathbb Q$ by its four coordinate equations.
Blockwise application gives a monic degree $2n$ rational annihilator in
the same way. These are carefully specified instances of the article's
noncommutative generalizations.

\subsection{What the octonion assertion does and does not give}

For one rational octonion $q=a_0+\sum_{i=1}^7a_ie_i$, direct use of the
multiplication table gives
$q\overline q=\sum_{i=0}^7a_i^2$ and $\overline q=2a_0-q$. Distributing
the product, which involves only two factors, proves
\[
 q^2-2a_0q+\left(\sum_{i=0}^7a_i^2\right)1=0.
\]
This is an elementary quadratic identity. It is not a proof of the
modified theorem for matrices over octonions mentioned in the article.
The latter requires a specified nonassociative evaluation convention
and a different theorem, and is not claimed here.

\subsection{Arithmetic status of these applications}

The constructions above use only fixed primitive recursive ring
operations, coefficient lists, products, bounded searches, and
primitive recursion with displayed bounds or parameters. Their stated
invariants are finite arrays of equalities and decidable inequalities.
Finite conjunctions are represented by a Boolean fold, so induction on
such invariants is quantifier-free induction after expansion of the
primitive recursive definitions. Algebraic structure axioms are used
as supplied equations; field inverses, root differences, generator
coordinates, and analytic tail estimates are additional data precisely
where indicated. This gives an explicit route to iPRA proofs of these
arithmetized statements. It does not turn an abstract host-language
proof about all rings, all modules, or all real functions into an
object-language iPRA derivation without the required coding and proof
translation.

\section{Lean formalization and the exact ProveIt interface}
\label{sec:lean}

\subsection{What was inspected}

The source audit is pinned to ProveIt commit
\texttt{73e6a339db7c72c6b13259059a8345b24385f6e3}.
The inspected project specifies Lean and Mathlib \texttt{v4.32.0}.
The relevant source root is
\begin{center}
\small\texttt{Logic/PrimitiveRecursiveArithmetic/ConciseIPRA/Lean/ConciseIPRA/}.
\end{center}
Its \texttt{PR.lean} defines the actual object theory, and
\texttt{NativeBiInterpretation.lean} provides the translation to the
native scalar calculus~\cite{proveit-pr,proveit-native}.
The packet inventory at
\path{Logic/PrimitiveRecursiveArithmetic/docs/proofs/README.md}
expressly distinguishes delivered research packets from rechecked
proofs~\cite{proveit-packets}. No theorem from those incoming packets is
used as a verified dependency here. The mathematical development above
is self-contained.

Inside the namespace \texttt{LeanProofs.ConciseIPRA}, the function-symbol
syntax is exactly
\begin{lstlisting}
inductive PR : Nat -> Type
  | zero : PR 0
  | succ : PR 1
  | proj {k : Nat} (i : Fin k) : PR k
  | comp {m k : Nat} (f : PR m)
      (gs : Fin m -> PR k) : PR k
  | prec {k : Nat} (f : PR k)
      (g : PR (k + 2)) : PR (k + 1)
\end{lstlisting}
The recursion argument comes first. Terms are \texttt{PTm n}, formulas
are \texttt{Fm PTm n}, the axiom family is \texttt{PRAAx}, and an internal
derivation has type \texttt{Prf PRAAx Gamma phi}. The wrapper for
quantifier-free induction is \texttt{PRA.qfInd}; it requires an explicit
proof that the induction formula satisfies \texttt{Fm.QF}.

\subsection{The three delivered Lean files}

\texttt{NatKernelSpec.lean} gives concrete natural-number definitions of
the coding and checking algorithm. Its functions include
\texttt{cons}, \texttt{head}, \texttt{tail}, \texttt{lookup},
\texttt{zadd}, \texttt{zmul}, dense polynomial addition and convolution,
matrix multiplication, permutation signs, determinants, minors,
adjugates, characteristic coefficients, and finally
\texttt{cayleyResidual} and \texttt{chCheck}. It imports only
\texttt{Std}. For the characteristic polynomial it stores $t$ as one
extra generic variable, then extracts its coefficients. This is
equivalent to the separate $t$-coefficient list used in the manuscript.
Its fixed-cap convolution is multiplication in a truncated polynomial
ring; the support proof in Section~\ref{sec:foundations} is essential
to showing that no relevant coefficient was lost.

These are host-level definitions of numerical functions. The remaining
syntactic step is to construct their \emph{particular} codes in
\texttt{PR k}, with the representation and defining-equation proofs.
The fixed functions passed to folds are implementation abbreviations,
not object-language quantification over all functions.

\texttt{ObjectLanguageTargets.lean} defines the exact target formula.
In schematic ASCII notation, its central definition is
\begin{lstlisting}
def checkStatement (check : PR 1) : Fm PTm 0 :=
  Fm.all (Fm.eq
    (PTm.app check (fun _ : Fin 1 => PTm.var 0))
    (PTm.succ PTm.zero))

def MainDerivationTarget (check : PR 1) : Prop :=
  Prf PRAAx [] (checkStatement check)
\end{lstlisting}
The argument must be the actual code \texttt{chCheckCode : PR 1}
representing the specified numerical checker. The file supplies no
postulated inhabitant of \texttt{MainDerivationTarget chCheckCode}.
Once an inhabitant is constructed, the existing theorem
\texttt{pra\_iff\_native} transfers it to
\texttt{NPrf [] (NFm.ofPRA (checkStatement chCheckCode))}.
That equivalence does not turn arbitrary Lean truth into iPRA provability.

\texttt{HostReference.lean} records ordinary Mathlib comparison targets.
For example,
\begin{lstlisting}
theorem cayleyHamilton (n : Nat)
    (A : Matrix (Fin n) (Fin n) R) :
    aeval A A.charpoly = 0 :=
  Matrix.aeval_self_charpoly A
\end{lstlisting}
Mathlib's theorem and its monic-remainder corollaries are useful for
checking the intended host-level mathematical statement~\cite{mathlib}.
They do not discharge the internal proof task. The three files contain
no added \texttt{axiom} declarations or \texttt{sorry} proof bodies.
Their absence does not mean that the missing target theorem has been
proved: it is deliberately left as a specification.

\subsection{A concrete route to the complete internal proof}

The following order avoids a hidden reflection theorem or an accidental
use of quantified induction.
\begin{enumerate}
\item \textbf{Coding and arithmetic.} Assemble \texttt{PR} syntax for
the finitely many list, signed-integer, comparison, iteration, and
bounded-sum helpers. Derive their equations with
\texttt{PRA.proj}, \texttt{PRA.comp}, \texttt{PRA.precZero}, and
\texttt{PRA.precSucc}. Prove equality tests, finite-sum reindexing, and
two-element cancellation internally.
\item \textbf{Polynomial arrays.} Prove convolution laws coefficient by
coefficient, packed into Boolean equalities. Prove support preservation,
padding compatibility, and specialization of normalized polynomials.
These lemmas give the precise link between the dense checker and the
ordinary integer polynomial notation.
\item \textbf{Matrices and permutations.} Prove the finite matrix ring
laws; implement insertion/deletion of a row and column, transpositions,
and the sign parity calculations. Derive cofactor expansion by the
explicit permutation bijection, then the two adjugate identities.
\item \textbf{Characteristic coefficients.} Construct the finite
coefficient matrices $B_j$. Prove $c_n=1$, $B_{n-1}=I$, and the
coefficient relations used in Theorem~\ref{thm:CH}. These are equalities
of coded arrays, not an application of an evaluation homomorphism whose
commutation hypotheses have not been proved.
\item \textbf{Telescoping.} Apply \texttt{PRA.qfInd} to the packed
version of
$\sum_{j=0}^{k}c_jX^j=-X^{k+1}B_k$ for the bounded prefix.
Handle the final leading coefficient and the empty-matrix branch.
Convert equality of the resulting array to \texttt{chCheck(n)=1},
universally close $n$, and obtain the actual
\texttt{Prf PRAAx [] (checkStatement chCheckCode)}.
\item \textbf{Applications.} Add the explicit monic-division,
Newton/Bell, Gaussian-elimination, interpolation, and finite-tail
invariants already stated in the respective proofs. Keep the input
certificates as assumptions until discharged for a particular
presentation. Apply the native translation only after each internal
derivation has been constructed.
\end{enumerate}

The accompanying \texttt{FORMALIZATION\_PLAN.md} supplies a more detailed
module table, coding equations, finite bounds, and source hashes. This
is substantive remaining formalization work, mainly the internal
finite-array and determinant library. A proof assistant cannot replace
it by observing that the answer of the executable checker is true.

\subsection{Verification status}

No Lean or Lake executable was available in this execution environment.
Consequently neither these outline files nor the upstream project was
compiled here. Their import paths and principal API names were checked
against the pinned source. The delivered PDF was compiled from the
delivered LaTeX, and the executable Python checks were run.

The core verification program computes signed permutation determinants
independently of a characteristic-polynomial library. It checks generic
matrices in dimensions $0$ through $4$ by exact integer-polynomial
normalization, $56$ integer matrices, $180$ matrix instances modulo
$2,4,6,9$, and $30$ rational quaternion representation examples.
It checks both adjugate identities, all coefficient recurrences,
Cayley--Hamilton, polynomial adjugates, and power reduction where
specified. It also verifies the displayed counterexamples and finite
factorial-tail sample bounds.

The separate trace program checks $43$ integer matrices in dimensions
$0$ through $4$ for Newton, Bell, Hessenberg, exact integer
Faddeev--LeVerrier, and truncated exponential/logarithm identities,
as well as the trace obstructions in characteristics $2,3,5$.
These finite checks can detect transcription and implementation errors;
the all-dimension results rest on the human proofs and their induction
accounting, not on extrapolation from the checks.



\section*{Source notes}
\addcontentsline{toc}{section}{Source notes}
The online article sets the scope, rather than serving as a collection
of unchecked axioms. The determinant and module arguments are classical
finite algebra, with related general formulations in
\cite{stacks,conrad}. Stein's account provides context for the
algebraic-number application~\cite{stein}. The quaternion discussion
uses explicit coordinate representations; Zhang's paper is relevant
background~\cite{zhang}. Tian's octonion-matrix work is cited only to
identify the more general result that is \emph{not} proved here~\cite{tian}.
The repository and Mathlib citations document actual interfaces, not
a claim that those interfaces automatically internalize these proofs.

\begin{thebibliography}{99}
\bibitem{wikipedia}
Wikipedia contributors, \emph{Cayley--Hamilton theorem}, inspected
revision \texttt{1375228092}, last edited 16 September 2026.
\url{https://en.wikipedia.org/w/index.php?title=Cayley%E2%80%93Hamilton_theorem&oldid=1375228092}.

\bibitem{stacks}
The Stacks Project Authors, \emph{Cayley--Hamilton}, Section 10.16,
Tag 05G6, especially Lemmas 10.16.1--10.16.3.
\url{https://stacks.math.columbia.edu/tag/05G6}.

\bibitem{conrad}
Brian Conrad, \emph{Generalized Cayley--Hamilton and integrality},
Math 210B notes, Stanford University.
\url{https://math.stanford.edu/~conrad/210BPage/handouts/cayleyhamilton.pdf}.

\bibitem{stein}
William A. Stein, \emph{Algebraic Number Theory, a Computational Approach},
Chapter 2, Basic Commutative Algebra.
\url{https://wstein.org/books/ant/ant/Ch2.html}.

\bibitem{zhang}
Fuzhen Zhang, Quaternions and matrices of quaternions,
\emph{Linear Algebra and its Applications} \textbf{251} (1997), 21--57.
DOI: \href{https://doi.org/10.1016/0024-3795(95)00543-9}{10.1016/0024-3795(95)00543-9}.
Author institution record:
\url{https://nsuworks.nova.edu/math_facarticles/100/}.

\bibitem{tian}
Yongge Tian, \emph{Matrix representations of octonions and their
applications}, arXiv:math/0003166v2 (2000).
\url{https://arxiv.org/abs/math/0003166v2}.

\bibitem{proveit-pr}
Vladimir Reshetnikov, ProveIt, \texttt{ConciseIPRA/PR.lean},
commit \texttt{73e6a339db7c72c6b13259059a8345b24385f6e3}.
\url{https://github.com/VladimirReshetnikov/ProveIt/blob/73e6a339db7c72c6b13259059a8345b24385f6e3/Logic/PrimitiveRecursiveArithmetic/ConciseIPRA/Lean/ConciseIPRA/PR.lean}.

\bibitem{proveit-native}
Vladimir Reshetnikov, ProveIt, \texttt{NativeBiInterpretation.lean},
same pinned commit.
\url{https://github.com/VladimirReshetnikov/ProveIt/blob/73e6a339db7c72c6b13259059a8345b24385f6e3/Logic/PrimitiveRecursiveArithmetic/ConciseIPRA/Lean/ConciseIPRA/NativeBiInterpretation.lean}.

\bibitem{proveit-packets}
Vladimir Reshetnikov, ProveIt, \emph{Delivered iPRA proof packets},
same pinned commit.
\url{https://github.com/VladimirReshetnikov/ProveIt/blob/73e6a339db7c72c6b13259059a8345b24385f6e3/Logic/PrimitiveRecursiveArithmetic/docs/proofs/README.md}.

\bibitem{mathlib}
The Mathlib Community, \texttt{Mathlib.LinearAlgebra.Matrix.Charpoly.Basic}
and \texttt{Coeff}, inspected tag \texttt{v4.32.0}.
\url{https://github.com/leanprover-community/mathlib4/tree/v4.32.0/Mathlib/LinearAlgebra/Matrix/Charpoly}.
Documentation:
\url{https://leanprover-community.github.io/mathlib4_docs/Mathlib/LinearAlgebra/Matrix/Charpoly/Basic.html}.
\end{thebibliography}
\end{document}
